\documentclass[10pt]{article}
\usepackage[T1]{fontenc}
\usepackage{lmodern,amsmath,amssymb,amsthm,booktabs,array}
\usepackage[margin=1in]{geometry}
\usepackage{microtype}
\PassOptionsToPackage{hyphens}{url}
\usepackage[colorlinks=true,urlcolor=blue,citecolor=blue,linkcolor=blue,pdfencoding=auto]{hyperref}
\hypersetup{pdftitle={Leading Kazhdan--Lusztig coefficients with fully commutative lower endpoints in type E},
 pdfauthor={Tyson Gern}}
\newtheorem{theorem}{Theorem}[section]
\newtheorem{lemma}[theorem]{Lemma}
\newtheorem{corollary}[theorem]{Corollary}
\newtheorem{proposition}[theorem]{Proposition}
\theoremstyle{remark}
\newtheorem{remark}[theorem]{Remark}
\newcommand{\FC}{W_{\mathrm{FC}}}
\newcommand{\supp}{\operatorname{supp}}
\newcommand{\dpth}{\operatorname{dp}}
\title{Leading Kazhdan--Lusztig coefficients with
fully commutative lower endpoints in type $E$\thanks{Text licensed under
CC BY 4.0; accompanying code under the MIT License.}}
\author{Tyson Gern\thanks{Initial Capacity, tyson@initialcapacity.io,
ORCID 0009-0003-8288-8786}}
\date{October 2026}
\begin{document}
\maketitle
\begin{abstract}
We prove that the leading coefficient $\mu(x,w)$ of the Kazhdan--Lusztig
polynomial $P_{x,w}$ lies in $\{0,1\}$ whenever $x$ is fully commutative
and $W$ is the Weyl group of type $E_6$, $E_7$ or $E_8$. Together with the
known results in types $A$ and $D$, this settles for every finite simply
laced Coxeter group the question raised by Green in 2009 and called
Green's $0$--$1$ conjecture by Chmutov. The proof reduces, by Green's star
argument as adapted by Gern, to a classification of the terminal
($T$-avoiding) elements of $E_6$, $E_7$ and $E_8$, obtained by exhaustive
enumeration with parabolic pruning: the noncommuting terminals are Gern's
type-$D$ list together with exactly two new elements, of lengths $28$ in
$E_7$ and $50$ in $E_8$, and the coefficient statement then follows from
Fan's monomial cell theory, in Green's formulation, and parity. We also
construct infinite families of full-support terminal reflections in every
generalized type $E_{4r+1}$, $r\geq3$, and in $E_{10}$, with
$\mu(x,b)=0$ for every fully commutative $x$ and every member of the
family, and an affine $E_8$ family with exact lengths $27+92k$ and the
same vanishing. The bound fails in affine type $D_4$, where $\mu=2$
occurs, and it fails without the fully commutative hypothesis: $\mu=10$
occurs already in $E_6$.
\end{abstract}
\noindent{\small MSC 2020: 20F55 (primary), 20C08, 05E10.
Keywords: Kazhdan--Lusztig polynomials, fully commutative elements,
Coxeter groups of type $E$, $T$-avoiding elements.}

\section{Introduction}\label{sec:intro}
Let $(W,S)$ be a Coxeter system with length function $\ell$ and Bruhat
order $\leq$, and let $P_{x,w}(q)$ be its ordinary equal-parameter
Kazhdan--Lusztig polynomials \cite{KL}. Put
$\mu(x,w)=[q^{(\ell(w)-\ell(x)-1)/2}]P_{x,w}(q)$ if $x<w$ and
$\ell(w)-\ell(x)$ is odd, and $\mu(x,w)=0$ otherwise. A \emph{cover} is a
pair $x<w$ with $\ell(w)-\ell(x)=1$; then $P_{x,w}=1$ and $\mu(x,w)=1$.
Write $L(w)=\{s\in S:\ell(sw)<\ell(w)\}$ and
$R(w)=\{s\in S:\ell(ws)<\ell(w)\}$ for the descent sets and $\supp(w)$
for the set of generators in a reduced word of $w$; $w$ has \emph{full
support} if $\supp(w)=S$. Following Stembridge \cite{Stembridge}, $w$ is
\emph{fully commutative} (FC) if any two of its reduced words are related
by interchanging adjacent commuting generators; in a simply laced system
this holds if and only if no reduced word of $w$ contains a factor $sts$
with $m(s,t)=3$ \cite[Proposition~2.1]{Stembridge}. The FC elements form
the set $\FC$. A set $I\subseteq S$ of pairwise commuting generators is an
\emph{independent set} of the Coxeter graph; $i(I)=\prod_{s\in I}s$ is a
\emph{commuting product}. An element is \emph{terminal} if no reduced word
for it begins or ends with two noncommuting generators, and a
\emph{noncommuting terminal} if moreover it is not a commuting product;
these are Gern's weakly bad and bad elements \cite[Definition~2.1.1]{Gern}
and Laird's $T$-avoiding and $T_2$-avoiding elements \cite{Laird}.

\begin{theorem}[Finite exceptional types]\label{thm:main}
If $W$ has type $E_6$, $E_7$ or $E_8$, then $\mu(x,w)\in\{0,1\}$ for all
$x\in\FC$ and $w\in W$.
\end{theorem}

\begin{corollary}[Finite simply laced groups]\label{cor:finiteADE}
If $W$ is a finite simply laced Coxeter group, $x\in\FC$ and $w\in W$,
then $\mu(x,w)\in\{0,1\}$.
\end{corollary}

Our second theme is arbitrarily large rank. Throughout, the generalized
diagram $E_n$, $n\geq6$, is the chain $0\text{--}1\text{--}\cdots\text{--}(n-2)$ with node $n-1$
attached to node~$2$; thus $E_9=\widetilde E_8$, and $E_n$ is indefinite
for $n\geq10$. Section~\ref{sec:uniform} defines, for every $r\geq3$,
explicit real roots $\beta_{r,k}$ ($k\geq0$) of $E_{4r+1}$, with
reflections $b_{r,k}=r_{\beta_{r,k}}$, and an independent set $I_r$ of
size $2r+1$.

\begin{theorem}[A rank-uniform family]\label{thm:uniform}
For every $r\geq3$ the reflections $b_{r,k}$, $k\geq0$, are pairwise
distinct full-support terminal elements of $W(E_{4r+1})$ that are not fully
commutative, with $L(b_{r,k})=R(b_{r,k})=I_r$. For every fully commutative
$x$ and every $k\geq0$, $\mu(x,b_{r,k})=0$.
\end{theorem}

Theorem~\ref{thm:affine} gives an analogous family in $\widetilde E_8$
with exact lengths $27+92k$, and Proposition~\ref{prop:e10} one in
$E_{10}$; so the terminal obstructions to length-decreasing star reduction
cannot be exhausted by a finite list in $E_n$ for $n\geq9$.

Green \cite[Introduction]{GreenLeading} asked whether $\mu(x,w)\in\{0,1\}$
for fully commutative $x$ in a large class of Coxeter groups, ``possibly
all of them''; Chmutov \cite[abstract and \S2.3]{Chmutov} calls this
\emph{Green's $0$--$1$ conjecture}. It is known in types $A$ and
$\widetilde A$ (Green \cite[Theorem~4.1; Theorem~3.1 of the arXiv
version]{GreenLeading}), $B=C$ (Chmutov \cite[Theorem~2.3.8]{Chmutov}) and
$D$ (Gern \cite[Theorem~4.5.11]{Gern}); Theorem~\ref{thm:main} completes
the finite simply laced case, $F_4$, $H_3$ and $H_4$ remain open (dihedral
groups, $G_2$ included, have all $P_{x,w}=1$), and finiteness cannot be
dropped from
Corollary~\ref{cor:finiteADE}, since $\mu=2$ occurs in $\widetilde D_4$
(Section~\ref{sec:complements}).

Jones \cite[Theorems~1.2 and~6.1]{Jones}
proved the bound for a Deodhar upper endpoint in every finite Weyl group
($E_6$, $E_7$, $E_8$, $F_4$ and $G_2$ by computer), the only earlier
exceptional $0$--$1$ result, with the opposite endpoint hypothesis; when
both endpoints are fully commutative the bound holds in every $E_n$ by
Graham's thesis, as recalled in \cite[Remark~9.5]{GreenEn} and
\cite[Remark~7.12]{GreenJones}. The reduction to terminal upper endpoints
(Lemma~\ref{lem:reduction}) is Green's argument as adapted by Gern, and
the lower-endpoint lemma (Lemma~\ref{lem:maxI}) is Fan's monomial cell
theorem \cite{Fan} in Green's formulation \cite{GreenEn}; what is new is
the classification of the terminal elements of $E_6$, $E_7$ and $E_8$
(Proposition~\ref{prop:classification}) by exhaustive enumeration, with
parabolic pruning in place of the $696{,}729{,}600$ elements of $E_8$. The
noncommuting terminals turn out to be Gern's type-$D$ list together with
exactly two new elements, $w_7\in E_7$ of length $28$ and $w_8\in E_8$ of
length $50$ (Remark~\ref{rem:gern}); each has a unique fully commutative
lower endpoint satisfying both descent inclusions, with even length gap
for the new ones, so the only Kazhdan--Lusztig value needed is Gern's $D_6$
value $\mu(x_6,w_6)=1$ \cite[Lemma~4.5.5]{Gern} (du Cloux's
\textsf{Coxeter}; recomputed in Appendix~\ref{app:certificates}).

Infinite families of $T$-avoiding elements inside a fixed infinite Coxeter group
were known in affine type $A$ (Fan--Green
\cite[Proposition~3.1.2]{FanGreen}, fully commutative), affine type $C$
(Ernst \cite{Ernst}; Laird \cite[Theorem~4.1.5]{Laird}) and type $F_n$
(non-FC; Cross, Ernst, Hills-Kimball and Quaranta (2012) and
Gilbertson--Ernst (2013), as reported in \cite[\S3.3]{Laird}); Laird
\cite[\S5.3]{Laird} lists the affine type-$E$ questions as open. By Green's
theorem \cite[Theorem~6.3]{GreenStar} every $E_n$ is star reducible, so
every fully commutative terminal in $E_n$ is a commuting product, which is
why all families here are non-FC.

Exceptional Kazhdan--Lusztig computations feature in Geck's \emph{PyCox} \cite{GeckPyCox} and in
Geck--Halls's analysis of the cells of $E_8$ \cite{GeckHalls}, which
concern cell partitions and character values rather than every $\mu(x,w)$
with FC lower endpoint; fully commutative cell closure fails in type $E$
(Green--Losonczy \cite[Example~2.2.5]{GreenLosonczy}, credited to Graham
\cite[\S9.9]{Graham}), and a $W$-graph records $\mu(x,w)$ only through a
generator in $L(x)\setminus L(w)$ \cite[Remark~2.1]{GeckHalls}, whereas
our residual terminal pairs have equal descent sets. To our knowledge, the
exceptional bound, the terminal lists of Table~\ref{tab:terminals} and the
infinite families have not appeared before.

Sections~\ref{sec:reductions},
\ref{sec:finite} and~\ref{sec:families} contain the structural lemmas,
the finite theorem and the infinite families, Section~\ref{sec:complements}
complements and open problems, and Appendix~\ref{app:certificates} the
computational certificates.

\section{Structural reductions}\label{sec:reductions}

\subsection{Star reduction}
We use the elementary properties of $\mu$ from \cite[\S2.3 and
Theorem~4.2]{KL}, restated in \cite[Proposition~1.2.4]{Gern}.
\begin{lemma}[Descent criterion]\label{lem:descent}
Let $x<w$. If $s\in L(w)\setminus L(x)$, then $\mu(x,w)=0$ unless $x=sw$;
if $s\in R(w)\setminus R(x)$, then $\mu(x,w)=0$ unless $x=ws$.
\end{lemma}
\begin{proof}
This is \cite[(2.3.e) and (2.3.f)]{KL}; see
\cite[Proposition~1.2.4(3),(4)]{Gern}.
\end{proof}
For noncommuting generators $s,t$ with $m(s,t)=3$ let
$D_L(s,t)=\{w:|L(w)\cap\{s,t\}|=1\}$. For $w\in D_L(s,t)$ exactly one of
$sw,tw$ lies in $D_L(s,t)$; it is the left star ${}^*w$. Right stars
$w^*$ are defined symmetrically. Stars preserve full commutativity in
simply laced systems \cite[Proposition~2.10]{Shi}, \cite[Proposition~1.1.29]{Gern}.
Star transport \cite[Theorem~4.2]{KL} says that for $x,w\in D_L(s,t)$,
\begin{equation}\label{eq:star}
 \widetilde\mu(x,w)=\widetilde\mu({}^*x,{}^*w),
\end{equation}
where $\widetilde\mu(y,z)$ equals $\mu(y,z)$ if $y<z$, $\mu(z,y)$ if
$z<y$, and $0$ otherwise; the same holds for right stars.

\begin{lemma}[Terminal reduction]\label{lem:reduction}
Let $W$ be simply laced. If $\mu(x,b)\in\{0,1\}$ for every terminal upper
endpoint $b$ and every fully commutative $x$, then $\mu(x,w)\in\{0,1\}$
for all $x\in\FC$ and all $w\in W$.
\end{lemma}
\begin{proof}
This is the argument of Green \cite[proof of Theorem~3.1 (published
Theorem~4.1)]{GreenLeading}, as adapted by Gern \cite[Lemma~2.3.10 and
Theorem~4.5.11]{Gern}. We induct on $\ell(w)$. Put $d=\ell(w)-\ell(x)$;
only $x<w$ with $d\geq3$ odd needs proof. If $L(w)\not\subseteq L(x)$ or
$R(w)\not\subseteq R(x)$, Lemma~\ref{lem:descent} gives zero; so assume
both inclusions, whence $L(w)$ and $R(w)$ are commuting, like the descent
sets of the fully commutative $x$. If $w$ is terminal, the hypothesis
applies. Otherwise let $w=stv$ be reduced with $m(s,t)=3$ (a noncommuting
suffix is treated with right stars). Then $s\in L(w)$, $t\notin L(w)$
because $L(w)$ is commuting, and $s,t\notin L(v)$ ($t$ by reducedness, and
$s\in L(v)$ would give $w=stsv'=tstv'$ and $t\in L(w)$); so $\{tv,stv\}$
is the left $\{s,t\}$-string in $D_L(s,t)$ and ${}^*w=sw$ has length
$\ell(w)-1$. Descent inclusion gives $s\in L(x)$ and full commutativity
$t\notin L(x)$, so $x\in D_L(s,t)$ and ${}^*x$ is defined and fully
commutative. By \eqref{eq:star}, $\mu(x,w)=\widetilde\mu({}^*x,{}^*w)$; as
$\ell({}^*x)=\ell(x)\pm1$, the new length difference is $d$ or $d-2$, both
positive, so the value is zero if the new endpoints are incomparable and
lies in $\{0,1\}$ by induction otherwise.
\end{proof}

\subsection{Maximum commuting descents}
The monomial cell structure of the Temperley--Lieb quotient in type $E_n$
is due to Fan \cite{Fan}; we use Green's formulation \cite[\S4]{GreenEn}
(statement numbers follow arXiv:0704.0283v1). The next lemma makes
explicit the singleton-cell argument ending Green's proof of
\cite[Lemma~5.2]{GreenEn}; its type-$D$ case is \cite[Lemma~4.5.8]{Gern}.
\begin{lemma}[Maximum common commuting descents]\label{lem:maxI}
Let $n\geq6$, $W=W(E_n)$, and let $x\in\FC$ with $\supp(x)\subseteq K$.
Let $I$ be an independent set of maximum cardinality in the subgraph
induced by $K$. If $I\subseteq L(x)\cap R(x)$, then $x=i(I)$.
\end{lemma}
\begin{proof}
Let $a_{\rm TL}(x)$ be the maximum size of a commuting product occurring as
a length-additive factor of $x$ \cite[Definition~4.1]{GreenEn}, Fan's
$a$-function \cite[Definition~2.3.1, p.~141]{Fan}. The descent inclusions
give a prefix and a suffix $i(I)$, so $a_{\rm TL}(x)\geq|I|$, while support
and maximality give $a_{\rm TL}(x)\leq|I|$. By
\cite[Proposition~4.4(iv)]{GreenEn}, which is
\cite[Proposition~4.4.3, p.~148]{Fan}, and its left-hand version (Fan's
dual statement, p.~148), $x$ lies in both the monomial right cell and the
monomial left cell of $i(I)$, and by \cite[Proposition~4.4(v)]{GreenEn},
that is \cite[p.~140 and proof of Theorem~6.1.2, p.~154]{Fan}, this
intersection is $\{i(I)\}$. Fan proves these statements for the whole
series $E_n$, $n\geq6$; no identification of monomial cells with
Kazhdan--Lusztig cells is used.
\end{proof}

\begin{proposition}[Maximum-descent parity criterion]\label{prop:parity}
Let $n\geq6$ and let $b\in W(E_n)$ satisfy $L(b)=R(b)=I$, where $I$ is a
maximum independent set in the subgraph induced by $\supp(b)$. If
$\ell(b)\equiv|I|\pmod2$, then for every $x\in\FC$,
\[
 \mu(x,b)\in\{0,1\},\qquad
 x<b,\ \ell(b)-\ell(x)>1\ \Longrightarrow\ \mu(x,b)=0.
\]
\end{proposition}
\begin{proof}
The type-$D$ prototype is \cite[Lemmas~4.5.4 and 4.5.9]{Gern}. Let $x<b$
be fully commutative with $\ell(b)-\ell(x)>1$ and $\mu(x,b)\neq0$. By
Lemma~\ref{lem:descent}, $I\subseteq L(x)\cap R(x)$; the subword property
gives $\supp(x)\subseteq\supp(b)$, so Lemma~\ref{lem:maxI} with
$K=\supp(b)$ forces $x=i(I)$. Then $\ell(b)-\ell(x)=\ell(b)-|I|$ is even
and $\mu(x,b)=0$, a contradiction. Covers have $\mu=1$.
\end{proof}
\begin{remark}\label{rem:parityremark}
Proposition~\ref{prop:parity} does not require $b$ to be terminal. The
parity hypothesis cannot be dropped: Section~\ref{sec:affineextra} exhibits
elements of $\widetilde E_8$ satisfying every other hypothesis with
$\mu(i(I),b)=1$.
\end{remark}
\begin{remark}\label{rem:indep}
The independence number of $E_n$ is $\lceil n/2\rceil$: excluding node
$2$ leaves three paths with maximum total $2+\lceil(n-4)/2\rceil$, and
including it gives $2+\lceil(n-5)/2\rceil$, the larger of the two being
$\lceil n/2\rceil$.
\end{remark}

\subsection{Geometric tests and covers of terminal reflections}\label{sec:covers}
Normalize simple roots by $(\alpha_j,\alpha_j)=2$ and
$(\alpha_s,\alpha_t)=-1$ for adjacent nodes, and write $s\sim t$ when $s$
and $t$ are adjacent, i.e.\ $m(s,t)=3$. For a real root $\beta$ the
reflection $r_\beta$ acts by $r_\beta(\rho)=\rho-(\rho,\beta)\beta$; $W$
acts on the left and the matrix of $w$ has columns $w(\alpha_j)$. Root
inequalities refer to the simple-root basis, and $s\in R(w)$ if and only
if $w(\alpha_s)<0$. Since $(ws)(\alpha_t)=w(\alpha_s+\alpha_t)$ for
$s\sim t$, the right-terminal test is
\begin{equation}\label{eq:terminaltest}
 s\in R(w)\quad\Longrightarrow\quad
 w(\alpha_s+\alpha_t)>0\quad\text{for every }t\sim s,
\end{equation}
and testing $w^{-1}$ gives the left-terminal condition.
The \emph{depth} $\dpth(\beta)$ of a positive root $\beta$ is the least $d$
such that $\beta=s_1\cdots s_{d-1}\alpha_t$ for simple reflections $s_i$
and a simple root $\alpha_t$ (cf.\ \cite[\S4.6]{BB}).
\begin{lemma}[Reflections]\label{lem:reflections}
Let $(W,S)$ be simply laced, let $\beta$ be a positive root and
$d=\dpth(\beta)$.
\begin{enumerate}
\item[(i)] $\ell(r_\beta)=2d-1$, and for every expression
$\beta=s_1\cdots s_{d-1}\alpha_t$ of minimal length the palindromic word
$s_1\cdots s_{d-1}\,t\,s_{d-1}\cdots s_1$ is a reduced word for $r_\beta$.
\item[(ii)] If $s\in S$ and $(\beta,\alpha_s)<0$, then $s\beta$ is a
positive root of depth $d+1$, $\ell(sr_\beta)=\ell(r_\beta)+1$ and
$\ell(sr_\beta s)=\ell(r_\beta)+2$. If $(\beta,\alpha_s)>0$ and
$\beta\neq\alpha_s$, then $s\beta$ is a positive root of depth $d-1$.
\item[(iii)] $r_\beta$ is fully commutative if and only if $\beta$ is
simple.
\end{enumerate}
\end{lemma}
\begin{proof}
Recall that $\ell(ws)>\ell(w)$ if and only if $w(\alpha_s)>0$, and that
$s$ negates only $\pm\alpha_s$. Let $(\beta,\alpha_s)=-c<0$. Then
$s\beta=\beta+c\alpha_s>0$; $r_\beta(\alpha_s)=\alpha_s+c\beta>0$ gives
$s\notin R(r_\beta)=L(r_\beta)$, so $\ell(sr_\beta)=\ell(r_\beta)+1$; and
$(sr_\beta)(\alpha_s)=-\alpha_s+c\,s\beta=c\beta+(c^2-1)\alpha_s>0$ gives
$\ell(sr_\beta s)=\ell(r_\beta)+2$. As $sr_\beta s=r_{s\beta}$, this shows
$\ell(r_{s\beta})=\ell(r_\beta)+2$ when $(\beta,\alpha_s)<0$ and, applied
to $s\beta$, $\ell(r_{s\beta})=\ell(r_\beta)-2$ when $(\beta,\alpha_s)>0$,
$\beta\neq\alpha_s$. Now $\ell(r_\beta)=2\dpth(\beta)-1$ by induction on
the odd number $\ell(r_\beta)$ (trivial for $\ell(r_\beta)=1$): if
$\ell(r_\beta)\geq3$, some $s$ has
$(\beta,\alpha_s)>0$, $\beta\neq\alpha_s$ (as $(\beta,\beta)=2$,
$\beta\geq0$), so $\beta'=s\beta>0$ has
$\ell(r_{\beta'})=\ell(r_\beta)-2=2\dpth(\beta')-1$ and
$\dpth(\beta)\leq\dpth(\beta')+1$ gives $2\dpth(\beta)-1\leq\ell(r_\beta)$,
while the palindromic word gives $\ell(r_\beta)\leq2d-1$. This proves (i);
(ii) follows from (i) and the two length formulas. For (iii), if $d\geq2$
the reduced word of (i) contains $s_{d-1}\,t\,s_{d-1}$, which equals $t$
unless $m(s_{d-1},t)=3$; so it is a braid factor and $r_\beta\notin\FC$
\cite[Proposition~2.1]{Stembridge}.
\end{proof}

\begin{lemma}[Covers of terminal reflections]\label{lem:covers}
Let $n\geq6$, $W=W(E_n)$, and let $b=r_\beta$ for a positive root $\beta$
with at least three nonzero coordinates. Suppose $L(b)=R(b)=I$, where $I$
is a maximum independent set of the subgraph induced by $\supp(b)$ and
$|I|$ is odd. Then no Bruhat cover of $b$ is fully commutative.
Consequently $\mu(x,b)=0$ for every fully commutative $x$.
\end{lemma}
\begin{proof}
Let $x<b$ with $\ell(x)=\ell(b)-1$. If $I\subseteq L(x)\cap R(x)$ and
$x\in\FC$, then $\supp(x)\subseteq\supp(b)$ and Lemma~\ref{lem:maxI} gives
$x=i(I)$ of odd length $|I|$, impossible as $\ell(b)$ is odd. Otherwise
some $s\in I$ has $s\notin L(x)$ (or, symmetrically, $s\notin R(x)$); the
lifting property \cite[Proposition~2.2.7]{BB} gives $x\leq sb$, hence
$x=sb=(bs)^{-1}$ by length and involutivity, and it suffices to show
$bs\notin\FC$. Put
$m=(\alpha_s,\beta)$; $s\in R(b)$ gives $b(\alpha_s)=\alpha_s-m\beta<0$,
so $m\geq1$, and $s\beta=\beta-m\alpha_s$ is a positive non-simple root,
its coordinates off $s$ being those of $\beta$. Since $sbs=r_{s\beta}$ and
$(\alpha_s,s\beta)=-m<0$, Lemma~\ref{lem:reflections}(ii) gives
$\ell(s\,r_{s\beta})=\ell(r_{s\beta})+1$, so $s$ followed by the
palindromic word of Lemma~\ref{lem:reflections}(i) for $r_{s\beta}$ is a
reduced word for $bs$ containing a braid factor; hence $bs\notin\FC$.
Non-covers are handled by Proposition~\ref{prop:parity}, since $\ell(b)$
and $|I|$ are odd.
\end{proof}

\section{Finite exceptional types}\label{sec:finite}
\subsection{The terminal classification}\label{sec:classification}
A fully commutative $x$ is \emph{eligible} for a terminal $b$ if
$x\leq b$, $L(b)\subseteq L(x)$ and $R(b)\subseteq R(x)$; by
Lemma~\ref{lem:descent}, every ineligible non-cover has $\mu(x,b)=0$.
Elements are written as strings of node labels in the numbering of
Section~\ref{sec:intro}: the chain $0\text{--}1\text{--}\cdots\text{--}(n-2)$ with node $n-1$
attached to node~$2$.\footnote{In Bourbaki's labelling of $E_n$ (chain
$1\text{--}3\text{--}4\text{--}\cdots\text{--}n$, node $2$ attached to node $4$), our node $0$ is
Bourbaki's node $1$, our node $j$ with $1\leq j\leq n-2$ is Bourbaki's
node $j+2$, and our node $n-1$ is Bourbaki's node $2$.} A string denotes
the product of the corresponding simple reflections in the order written,
with the left action and column convention of Section~\ref{sec:covers}.
Every element displayed in Table~\ref{tab:terminals} and in the affine
family of Section~\ref{sec:affine} is an involution, so a string and its
reverse denote the same element.

\begin{proposition}[Computational classification]\label{prop:classification}
The noncommuting terminals in $E_6$, $E_7$ and $E_8$ are exactly the eleven
elements in Table~\ref{tab:terminals}. Both descent sets of each of them
equal the support of the commuting product $x$ shown in the table. The
remaining terminals are commuting products: $22$, $36$ and $58$ of them,
respectively, including the identity.
\end{proposition}
Extending Gern's notation \cite[Theorem~2.3.6]{Gern}, we write $w_4$ and
$w_6$ for his elements of lengths $7$ and $15$ (viewed inside $E_n$ as in
Remark~\ref{rem:gern}), and $w_7$, $w_8$ for the two new elements of
lengths $28$ and $50$, where $w_8$ is the concatenation of the two strings
\[
 (7534231270123456210321432)(5437210321432543721324357).
\]
\begin{table}[htbp]
\centering\small
\begin{tabular}{@{}clrrrl@{}}
\toprule
Type & Terminal $b$ & $\ell(b)$ & $x=i(I)$ & $\ell(b)-\ell(x)$ & Source\\
\midrule
$E_6$ & $1325213$ & 7 & $135$ & 4 & $w_4$\\
$E_7$ & $1326213$ & 7 & $136$ & 4 & $w_4$\\
$E_7$ & $13256213$ & 8 & $1356$ & 4 & $w_4s_5$\\
$E_7$ & $132543621324356$ & 15 & $1356$ & 11 & $w_6$\\
$E_7$ & $1325436210321432543621324356$ & 28 & $1356$ & 24 & $w_7$\\
$E_8$ & $1327213$ & 7 & $137$ & 4 & $w_4$\\
$E_8$ & $13257213$ & 8 & $1357$ & 4 & $w_4s_5$\\
$E_8$ & $61327213$ & 8 & $1367$ & 4 & $w_4s_6$\\
$E_8$ & $132543721324357$ & 15 & $1357$ & 11 & $w_6$\\
$E_8$ & $1325437210321432543721324357$ & 28 & $1357$ & 24 & $w_7$\\
$E_8$ & $w_8$ (displayed above) & 50 & $1357$ & 46 & $w_8$\\
\bottomrule
\end{tabular}
\caption{The noncommuting terminals of $E_6$, $E_7$, $E_8$: a reduced
word, the length, the unique eligible fully commutative lower endpoint
(Corollary~\ref{cor:eligible}), the gap, and the identification of
Remark~\ref{rem:gern}. Only the rows of length $15$, Gern's $D_6$ pair,
have odd gap.}
\label{tab:terminals}
\end{table}

\begin{corollary}\label{cor:eligible}
Each terminal in Table~\ref{tab:terminals} has the unique eligible fully
commutative lower endpoint $x$ displayed there.
\end{corollary}
\begin{proof}[Proof of Corollary~\ref{cor:eligible}]
The induced support graphs have types $D_4$, $D_4\sqcup A_1$, $D_6$,
$E_7$ and $E_8$, with independence numbers $3,4,4,4,4$: the matchings
$(1,2)$ in the rows of lengths $7$ and $8$, $(1,2),(3,4)$ in the rows of
length $15$, $(0,1),(2,3),(4,5)$ in the rows of length $28$ and
$(0,1),(3,4),(5,6),(2,7)$ in the row of length $50$ bound every independent
set by the number of support nodes minus the matching size, and the
displayed descent set attains the bound (compare Remark~\ref{rem:indep}).
If $y$ is eligible, then $\supp(y)\subseteq\supp(b)$ and
Lemma~\ref{lem:maxI} forces $y=x$; conversely $x$, a commuting product of
left descents of $b$, is a reduced prefix of $b$, hence below $b$ and
eligible.
\end{proof}

\begin{remark}[Gern's list plus two]\label{rem:gern}
The parabolic subgroup of $E_n$ on $\{1,n-1,2,\ldots,n-2\}$ has type
$D_{n-1}$; sending Gern's labels $1,2,j$ ($3\leq j\leq7$) to $1,n-1,j-1$
identifies it with Gern's $W(D_{n-1})$, and then the rows of
Table~\ref{tab:terminals} of lengths $7$, $8$ and $15$ are Gern's bad
elements $w_4$, $w_4u$ and $w_6$ of \cite[Theorem~2.3.6]{Gern}, while the
$E_8$ row of length $28$ is $w_7$ in the parabolic subgroup on
$\{0,\ldots,5,7\}$. The $E_n$ form of \cite[Theorem~2.3.6]{Gern} holds for
$n=6,7,8$: the noncommuting terminals of $E_n$ are exactly the $w_ku$ with
$k\in\{4,6,7,8\}$ and $w_k\in W(E_n)$ (so $k=4$ for $n=6$ and $k\leq n$ for
$n=7,8$), and $u$ a commuting product of generators neither in nor adjacent
to $\supp(w_k)$; for $n=7,8$, $w_n$ is the unique longest one (matrix
equality; exhaustive enumeration of $D_5$, $D_6$, $D_7$ for
the type-$D$ parabolics). Each $w_k$ is an involution with $L=R=I$: $w_4$
a reflection, $w_6$ a product of three orthogonal reflections, $w_7$ and
$w_8$ of involution types $4A_1$ and $2A_1$; the palindromic layered
words, the negated roots and the right weak order chain $w_4<w_6<w_7<w_8$
are in the certificate (Appendix~\ref{app:certificates}). For $b=w_7,w_8$,
$\mu(i(I),b)=0$ also follows without parity: $J=I\cup\{0\}$ has type
$A_2\times A_1^3$ and $\ell(w_0(J))=6$, and $b=u\,w_0(J)\,v$
length-additively with lengths $11+6+11$ and $11+6+33$ (words in the
certificate), so $b\leq_{LR}w_0(J)$ \cite[\S2.3]{KL},
$a(b)\geq a(w_0(J))=6$ for Lusztig's $a$-function
\cite[Proposition~13.8 and Conjecture~14.2 (P4), (P12), proved for finite
Weyl groups in \S15.1]{LusztigUP}, \cite[Lemmas~1.3.7 and 1.3.8]{Gern},
and, as $I\subseteq L(b)$, $P_{i(I),b}=P_{e,b}$ \cite[\S2.3]{KL} has
degree at most $(\ell(b)-a(b))/2\leq(\ell(b)-6)/2$
\cite[Proposition~1.3.5]{Gern}, citing \cite[1.3(a)]{LusztigII}, below the
degree $(\ell(b)-5)/2$ of $\mu(i(I),b)$. Likewise $a(w_6)\geq5$,
consistent with $a(w_6)=5$ \cite[Proposition~4.5.1]{Gern}.
\end{remark}

\subsection{Exhaustiveness, and the proofs of Theorem~\ref{thm:main} and
Corollary~\ref{cor:finiteADE}}\label{sec:exhaustive}
All computations use the faithful integral geometric representation: right
multiplication by $s_i$ replaces the columns $a_j=w(\alpha_j)$ by
$a_i\mapsto-a_i$ and $a_j\mapsto a_j+a_i$ for $j\sim i$; signs give
descents and \eqref{eq:terminaltest} tests terminality on both sides. In
$E_6$ and $E_7$, breadth-first closure under all generators enumerates the
whole group. For $E_8$, fix $J\subseteq S$ and write $w=av$
length-additively with $a\in W^J$ (minimal coset representatives) and
$v\in W_J$; call $w$ \emph{right-terminal} if no reduced word of $w$ ends
in a noncommuting pair.
\begin{lemma}[Parabolic pruning]\label{lem:pruning}
If $w=av$ is right-terminal, then $v$ is right-terminal in $W_J$.
\end{lemma}
\begin{proof}
A reduced word for $v$ ending in adjacent generators, prefixed by a
reduced word for $a$, gives such a word for $w$, by length additivity.
\end{proof}
So one enumerates the products $av$ with $v$ right-terminal, applies
\eqref{eq:terminaltest} in the ambient group and tests inverses. The
implementation recurses along
$E_8\supset E_7\supset E_6\supset D_5\supset D_4\supset A_3\supset
A_2\supset A_1\supset\{e\}$, obtaining the minimal coset representatives
at each step by closure under left multiplication and stripping of right
$J$-descents ($W$ acts transitively on $W/W_J$); Lemma~\ref{lem:pruning}
and induction give completeness at every stage. The final $E_7$ parabolic
has $576$ right-terminals and $240$ cosets, their products give $2{,}160$
right-terminals in $E_8$, and the inverse test leaves exactly the $64$
terminals of Proposition~\ref{prop:classification}. With
Corollary~\ref{cor:eligible}, the finite theorem thus needs neither an
enumeration of pairs with FC lower endpoint nor a table of
Kazhdan--Lusztig polynomials.

\begin{proof}[Proof of Theorem~\ref{thm:main}]
By Lemma~\ref{lem:reduction} it suffices to treat terminal $b$. If $b$ is
a commuting product, every $x<b$ is a proper subproduct, so some
$s\in L(b)\setminus L(x)$ exists and Lemma~\ref{lem:descent} gives
$\mu(x,b)\in\{0,1\}$. For a noncommuting terminal,
Corollary~\ref{cor:eligible} and Lemma~\ref{lem:descent} leave one
candidate non-cover pair $(x,b)$ per row of Table~\ref{tab:terminals}, and
parity disposes of all rows except the two of length $15$. These are the
same pair in a $D_6$ parabolic subgroup: in $E_n$, $n=7,8$, mapping Gern's
labels $1,\ldots,6$ to $1,n-1,2,3,4,5$ identifies it with Gern's
$x_6=(-1,-2,4,3,6,5)$ and $w_6=(-1,-6,3,-4,5,-2)$ in signed permutation
notation \cite[Example~1.1.6]{Gern}. Bruhat order and Kazhdan--Lusztig
polynomials of a standard parabolic subgroup agree with those of $W$, since
reduced words of $W_J$ are reduced in $W$ and the Hecke algebra of $W_J$ is
a bar-stable subalgebra containing its Kazhdan--Lusztig basis
\cite[Chapter~5]{BB}. Gern's Lemma~4.5.5 \cite{Gern}, a computation with du
Cloux's \textsf{Coxeter}, gives $\mu(x_6,w_6)=1$ (recomputed in
Appendix~\ref{app:certificates}). Hence every coefficient with terminal
upper endpoint lies in $\{0,1\}$.
\end{proof}

\begin{proof}[Proof of Corollary~\ref{cor:finiteADE}]
The irreducible finite simply laced types are $A$, $D$ and $E$; combine
Green's theorem, Gern's theorem and Theorem~\ref{thm:main}. For a direct
product $W=\prod_jW_j$, full commutativity is componentwise and
$P_{x,w}=\prod_jP_{x_j,w_j}$. Let $x<w$, let $d_j$ be the component length
differences, $d=\sum_jd_j$, and $m$ the number of $j$ with $d_j>0$. Even
$d$ gives $\mu=0$; for odd $d$ and $m\geq2$,
$\deg P_{x_j,w_j}\leq(d_j-1)/2$ gives $\deg P_{x,w}\leq(d-m)/2<(d-1)/2$,
so $\mu=0$; for $m=1$, $\mu(x,w)$ is the coefficient in the single
nontrivial component, which lies in $\{0,1\}$.
\end{proof}

\section{Infinite families}\label{sec:families}

\subsection{The rank-uniform construction}\label{sec:uniform}
Fix an integer $r\geq3$, put $n=4r+1$, and use the $E_n$ numbering of
Section~\ref{sec:intro}, so node $4r$ is attached to node $2$. Define
$\beta_r$ by its coordinates $(\beta_r)_0=r-1$, $(\beta_r)_1=r$,
$(\beta_r)_2=2r-1$, $(\beta_r)_{4r}=r$ and $(\beta_r)_j=2r-\lfloor j/2\rfloor$
for $3\leq j\leq4r-1$, where $(v)_j$ denotes the coefficient of $\alpha_j$
in $v$; so the long-arm coordinates are $2r-1,2r-2,2r-2,\ldots,1,1$. Put
$I_r=\{0,3,5,\ldots,4r-1,4r\}$ and $a=2r-4$. On the embedded affine $E_8$
subdiagram $\{0,\ldots,7,4r\}$ set
$(\delta_0,\ldots,\delta_7)=(2,4,6,5,4,3,2,1)$, $\delta_{4r}=3$,
$(\gamma_0,\ldots,\gamma_7)=(1,2,3,2,2,1,1,0)$, $\gamma_{4r}=1$, with all
other coordinates zero. Set $T=r_\gamma r_{\gamma+\delta}$ and, for
$k\geq0$,
\begin{equation}\label{eq:uniformroots}
 \beta_{r,k}=\beta_r-ak\gamma+(ak^2+rk)\delta,\qquad
 b_{r,k}=r_{\beta_{r,k}}.
\end{equation}

\begin{lemma}\label{lem:uniformseed}
For every $r\geq3$, $\beta_r$ is a real root of $E_{4r+1}$.
\end{lemma}
\begin{proof}
Deleting node $0$ leaves finite type $D_{4r}$; realize its simple roots in
an orthonormal basis $e_1,\ldots,e_{4r}$ as $\alpha_1=e_1+e_2$,
$\alpha_{4r}=e_2-e_1$, $\alpha_j=e_{j+1}-e_j$ ($2\leq j\leq4r-1$), take $z$
orthogonal to them with $(z,z)=2-r$, and put
$\alpha_0=z-\frac12\sum_ie_i$; these vectors have the Cartan inner
products of $E_{4r+1}$ and form a basis, the Cartan determinant $9-n$
being nonzero. Then $\beta_r=(r-1)z+\sum_ic_ie_i$ with $c_i=-(r-1)/2$ for
odd and $(3-r)/2$ for even $i$, and $(\beta_r,\alpha_0)=r-2$ gives
$s_0\beta_r=F\alpha_0$, where $F$ flips the $2r$ even coordinates. As
$F=\prod_{j=1}^r r_{e_{4j-2}-e_{4j}}r_{e_{4j-2}+e_{4j}}\in W(D_{4r})$
\cite[Example~1.1.6]{Gern}, $\beta_r=s_0F\alpha_0$ is a real root.
\end{proof}

\begin{lemma}\label{lem:rootaction}
Suppose $\gamma$ and $\gamma+\delta$ are real roots with
$(\gamma,\gamma)=2$ and $(\gamma,\delta)=(\delta,\delta)=0$, and let
$T=r_\gamma r_{\gamma+\delta}$. Then
\[
 (T-1)\rho=(\rho,\delta)\gamma-((\rho,\gamma)+(\rho,\delta))\delta,
 \qquad(T-1)^3=0,
\]
and if $(\beta,\delta)=-a$ and $(\beta,\gamma)=-c$, then
\begin{equation}\label{eq:rootiterate}
 T^k\beta=\beta-ak\gamma+(ak^2+ck)\delta\qquad(k\geq0).
\end{equation}
\end{lemma}
\begin{proof}
Compose the two reflection formulas. With $N=T-1$, $N\delta=0$ and
$N\gamma=-2\delta$ give $N^3=0$, and $N\beta=-a\gamma+(a+c)\delta$,
$N^2\beta=2a\delta$; expand $(1+N)^k$ binomially. The lemma holds on the
whole ambient root space, where $\delta$ need not be in the radical.
\end{proof}

\begin{lemma}\label{lem:uniformdescents}
For all $r\geq3$ and $k\geq0$, $\beta_{r,k}=T^k\beta_r$ is a real root with
every coordinate strictly positive. Its reflection $b_{r,k}$ is terminal
and has both descent sets equal to $I_r$.
\end{lemma}
\begin{proof}
The vector $\gamma$ is a root of the finite $E_8$ parabolic subsystem
obtained by omitting node $7$, and $\gamma+\delta$ is an affine real root
(lowering words to a simple root are recorded in the certificate,
Appendix~\ref{app:certificates}). As $(\beta_r,\gamma)=-r$ and
$(\beta_r,\delta)=-a$, \eqref{eq:rootiterate} gives
$T^k\beta_r=\beta_{r,k}$, a real root, and the binomial increments
$(a+r)\delta-a\gamma$ and $2a\delta$ are coordinatewise nonnegative for
$r\geq3$, so every coordinate of $\beta_{r,k}$ is strictly positive. The
pairings $m_j=(\alpha_j,\beta_{r,k})$ are
\begin{gather*}
 m_0=r-2,\quad m_1=2-r,\quad m_2=-1-ak,\quad m_8=-1-ak^2-rk,\quad
 m_{4r}=1+ak,\\
 m_j=(-1)^{j+1}(1+ak)\ \ (3\leq j\leq7),\qquad
 m_j=(-1)^{j+1}\ \ (9\leq j\leq4r-1),
\end{gather*}
positive exactly on $I_r$. Since $\beta_{r,k}>0$,
$b_{r,k}(\alpha_j)=\alpha_j-m_j\beta_{r,k}$ is negative exactly when
$m_j>0$, so $R(b_{r,k})=I_r=L(b_{r,k})$ by involutivity. For every edge
$s\sim t$ with $s\in I_r$ the sum $m_s+m_t$ is nonpositive (at $7-8$ it is
$-ak(k-1)-rk$, at $8-9$ it is $-ak^2-rk$, elsewhere by inspection), so
$b_{r,k}(\alpha_s+\alpha_t)=\alpha_s+\alpha_t-(m_s+m_t)\beta_{r,k}>0$,
which is \eqref{eq:terminaltest}; the left test follows by involutivity.
\end{proof}

\begin{proof}[Proof of Theorem~\ref{thm:uniform}]
The set $I_r$ is independent of size $2r+1$, and the matching
$(0,1)$, $(2,4r)$, $(3,4)$, $(5,6)$, \ldots, $(4r-3,4r-2)$ with $2r$ edges
leaves one vertex unmatched, so $I_r$ is a maximum independent set of
$E_{4r+1}$ (compare Remark~\ref{rem:indep}). By
Lemma~\ref{lem:uniformdescents}, $b_{r,k}$ is a terminal reflection with
$L=R=I_r$, $|I_r|$ is odd, and $\beta_{r,k}$ has $4r+1\geq3$ nonzero
coordinates, so Lemma~\ref{lem:covers} gives $\mu(x,b_{r,k})=0$ for every
fully commutative $x$, and $b_{r,k}\notin\FC$ by
Lemma~\ref{lem:reflections}(iii). Since $m_0=r-2>0$ and $\beta_{r,k}$ is
positive in every coordinate, column $0$ of $b_{r,k}$ changes every
coordinate row, whereas omitting a generator would leave its row
unchanged; so $b_{r,k}$ has full support. Distinct positive roots give
distinct reflections, and $(\beta_{r,k})_0=r-1+4k+(4r-8)k^2$ gives
distinctness.
\end{proof}

\begin{remark}\label{rem:uniformlengths}
At $r=3$ the seed is $\beta_3=(2,3,5,5,4,4,3,3,2,2,1,1,3)$ and
$\ell(b_{3,0})=75$, while the maximum lengths of fully commutative
elements of $E_{10},\ldots,E_{13}$ are $55$, $66$, $78$ and $92$
(Appendix~\ref{app:certificates}), so a length-based exclusion of fully
commutative covers already fails in rank $13$; the lengths
$\ell(b_{r,k})=8r^2+3+116k$ hold for every pair in the certificate ($r\leq8$,
$k\leq3$; the supplement's optional extended run adds $r=9$, $10$, $12$) but are not
needed. The restriction $r\geq3$ is structural: at $r=2$,
$\beta_2=\delta-\gamma$, $(\alpha_0,\beta_2)=0$, and the descent set has
even size $4$.
\end{remark}

\subsection{The affine family}\label{sec:affine}
Use $E_9$ numbering: chain $0\text{--}1\text{--}\cdots\text{--}7$ with node $8$ attached to $2$;
the affine node is $7$. Put
\[
 \delta'=(2,4,6,5,4,3,2,1,3),\qquad
 \gamma'=(1,1,1,1,1,0,0,0,0),\qquad I'=\{1,3,5,7,8\},
\]
and
\[
 c_0=312875645234123012856745231,
\]
a reduced word for an involution. Let $T'=r_{\gamma'}r_{\gamma'+\delta'}$,
so that $T'(\alpha)=\alpha-(\alpha,\gamma')\delta'$, and define
$c_k=T'^kc_0T'^{-k}$ for $k\geq0$.

\begin{theorem}[Explicit affine family]\label{thm:affine}
Every $c_k$ is a full-support terminal element of $W(\widetilde E_8)$ that
is not fully commutative, with
\[
 L(c_k)=R(c_k)=I',\qquad\ell(c_k)=27+92k.
\]
For every fully commutative $x\in W(\widetilde E_8)$ and every $k\geq0$,
$\mu(x,c_k)=0$.
\end{theorem}
\begin{proof}
The vector $\delta'$ is the null root and $\gamma'$ is a root of the
finite $E_8$ parabolic subsystem obtained by omitting node $7$.
Multiplying out the displayed word gives $c_0^2=1$. The columns
$c_0(\alpha_j)$ of the matrix of $c_0$ recorded in the affine certificate
(Appendix~\ref{app:certificates}; the certificate stores the matrix by rows)
are negative precisely at the indices
in $I'$, and for every edge incident to $I'$ the sum of the two columns is
a positive root; by \eqref{eq:terminaltest} and involutivity, $c_0$ is
terminal on both sides. Put $d=(2,-1,1,-1,1,-1,1,-1,-1)$. Then
$(\alpha_j,\gamma'-c_0\gamma')=d_j$, and conjugation gives
\begin{equation}\label{eq:affinecolumns}
 c_k(\alpha_j)=c_0(\alpha_j)+kd_j\delta'.
\end{equation}
As $d$ is negative on $I'$ and positive off $I'$, like the base columns,
both descent sets stay $I'$; on every edge $s\sim t$ with $s\in I'$ one
has $d_s+d_t\geq0$, so \eqref{eq:affinecolumns} preserves positivity of
$c_k(\alpha_s+\alpha_t)$, and $c_k$ is terminal. Full support holds at
$k=0$ by the displayed word and at $k\geq1$ because each descent column is
strictly negative in every coordinate, whereas omitting a generator would
leave its row unchanged. The $c_k$ are pairwise distinct by
\eqref{eq:affinecolumns}, since every $d_j\neq0$.

For the length, count affine inversions: for each of the $240$ finite
roots $\eta$ write $c_0\eta=\eta'+h\delta'$, let $\epsilon(\eta)$ be $0$
for positive and $1$ for negative $\eta$; the positive affine roots
$\eta+m\delta'$, $m\geq\epsilon(\eta)$, contribute
$\max(0,A_\eta+kB_\eta)$ to $\ell(c_k)$, where
$A_\eta=\epsilon(\eta')-h-\epsilon(\eta)$ and $B_\eta=-\sum_j\eta_jd_j$.
The nonzero contributions, grouped by $(A_\eta,B_\eta)$, are
\begin{center}
\begin{tabular}{@{}rrrrrr@{}}
\toprule
$(A,B)$ & $(0,1)$ & $(0,2)$ & $(1,0)$ & $(1,1)$ & $(1,2)$\\
Multiplicity & 44 & 8 & 1 & 20 & 6\\
\bottomrule
\end{tabular}
\end{center}
and the remaining $161$ vanish for all $k\geq0$; the sum is $27+92k$.

By Remark~\ref{rem:indep}, $I'$ is a maximum independent set, of size
$5$, and $\ell(c_k)$ is odd, so Proposition~\ref{prop:parity} gives
$\mu(x,c_k)=0$ for every fully commutative non-cover $x$. Let $x$ be a
fully commutative cover of $c_k$. If $I'\subseteq L(x)\cap R(x)$, then
Lemma~\ref{lem:maxI} gives $x=i(I')$ of length $5<\ell(c_k)-1$.
Otherwise, by the lifting property \cite[Proposition~2.2.7]{BB}, $x=sc_k$
or $x=c_ks$ for some $s\in I'$, as in the proof of
Lemma~\ref{lem:covers}. For $k\geq1$ these elements have length
$26+92k\geq118$, exceeding the maximum length $44$ of a fully commutative
element of $\widetilde E_8$ (complete catalogue of $44{,}199$ elements,
Appendix~\ref{app:certificates}, in agreement with \cite[\S5.2]{BJN}). For
$k=0$, every cover arises by deleting one letter from the displayed
reduced word, by the strong exchange condition \cite[Theorem~1.4.3]{BB};
the deletions give exactly $21$ distinct elements of length $26$, none
fully commutative (Appendix~\ref{app:certificates}). Hence $\mu(x,c_k)=0$ for every fully
commutative $x$. Finally, $c_k\notin\FC$, since otherwise
Lemma~\ref{lem:maxI} would force $c_k=i(I')$. (Lemma~\ref{lem:covers}
does not apply directly: $c_0$ is not a reflection.)
\end{proof}

\subsection{The \texorpdfstring{$E_{10}$}{E10} family}\label{sec:e10}
The first indefinite rank $E_{10}$ lies outside the sequence $E_{4r+1}$.
Embed the affine subdiagram on $\{0,\ldots,7,9\}$ and take
$\hat\delta=(2,4,6,5,4,3,2,1,0,3)$, $\hat\gamma=(1,2,3,2,2,1,1,0,0,1)$ (so
$\hat\gamma$ differs from $\gamma'$ of Section~\ref{sec:affine}),
$\hat\beta_0=(3,7,10,9,7,6,4,3,1,6)$, $\hat I=\{1,3,5,7,9\}$,
$\hat T=r_{\hat\gamma}r_{\hat\gamma+\hat\delta}$, and for $k\geq0$
\begin{equation}\label{eq:e10roots}
 \hat\beta_k=\hat\beta_0-k\hat\gamma+(k^2+3k)\hat\delta,\qquad
 \hat b_k=r_{\hat\beta_k}.
\end{equation}
\begin{proposition}\label{prop:e10}
The reflections $\hat b_k$, $k\geq0$, are pairwise distinct full-support
terminal elements of $W(E_{10})$ that are not fully commutative, with
$L(\hat b_k)=R(\hat b_k)=\hat I$. For every fully commutative $x$ and every
$k\geq0$, $\mu(x,\hat b_k)=0$.
\end{proposition}
\begin{proof}
A lowering word in the certificate (Appendix~\ref{app:certificates})
reduces $\hat\beta_0$ to $\alpha_9$; $\hat\gamma$, $\hat\gamma+\hat\delta$
are real roots of the embedded affine subsystem; and
$(\hat\beta_0,\hat\gamma)=-3$, $(\hat\beta_0,\hat\delta)=-1$, so
Lemma~\ref{lem:rootaction} gives $\hat\beta_k=\hat T^k\hat\beta_0$ with
nonnegative increments $4\hat\delta-\hat\gamma$, $2\hat\delta$, whence
$\hat\beta_k>0$ in every coordinate. The pairings
\[
 ((\alpha_j,\hat\beta_k))_{j=0}^9
 =(-1,\,1,\,-2-k,\,1+k,\,-1-k,\,1+k,\,-1-k,\,1+k,\,-1-k^2-3k,\,2+k)
\]
are positive exactly on $\hat I$ with nonpositive sums over adjacent
entries, giving both descent sets and terminality by
\eqref{eq:terminaltest}; $(\alpha_1,\hat\beta_k)=1$ gives full support
(column $1$ changes every row) and $(\hat\beta_k)_0=2k^2+5k+3$
distinctness. $\hat I$ is a maximum independent set of odd size $5$
(Remark~\ref{rem:indep}) and $\hat\beta_k$ has ten nonzero coordinates, so
Lemma~\ref{lem:covers} gives $\mu(x,\hat b_k)=0$ for every fully
commutative $x$, and $\hat b_k\notin\FC$ by Lemma~\ref{lem:reflections}(iii).
\end{proof}
Standard parabolic embeddings transport the $\widetilde E_8$ and $E_{10}$
families to larger $E_n$, with the supports of the smaller parabolic
subgroups.

\section{Complements and open problems}\label{sec:complements}

\subsection{Necessity of the fully commutative hypothesis}
A complete table of Kazhdan--Lusztig polynomials of $W(E_6)$ ($51{,}840$
elements; Appendix~\ref{app:certificates}) gives $\max\mu(x,w)=10$ over
all pairs, attained for instance by $x=01254012010$ ($\ell=11$, not FC)
and $w=123012523412301251234012$ ($\ell=24$, full support), with
\[
 P_{x,w}=1+13q+59q^2+121q^3+125q^4+61q^5+10q^6.
\]
Exactly $1{,}386$ pairs
have $\mu\geq2$ ($400$, $556$, $310$, $108$, $4$ and $8$ pairs with
$\mu=2,3,4,5,6,10$), all with non-FC lower endpoint, while all $6{,}431$
nonzero values with fully commutative lower endpoint equal $1$, a
brute-force confirmation of Theorem~\ref{thm:main} for $E_6$ (by parabolic
inclusion, $\mu\geq10$ occurs in $E_7$ and $E_8$ as well). Full
commutativity of the upper endpoint alone does not suffice either:
$x=5343010$ (not FC) and $w=0125342312501234$ (FC, $\ell=16$) have
\[
 P_{x,w}=1+8q+22q^2+20q^3+5q^4,\qquad \mu(x,w)=5.
\]

\subsection{The bound fails in affine type}
In $\widetilde D_4$, with $x$ the product of the four leaves and $c$ the
centre, $w=xcx$ has $P_{x,w}=1+3q+2q^2$ and $\mu(x,w)=2$; this star family
is covered by Mongelli \cite[Corollary~2 (arXiv v1, p.~10)]{Mongelli}, and
the value is recomputed in Appendix~\ref{app:certificates}. Hence
Corollary~\ref{cor:finiteADE} is sharp within simply laced type, and the
universal form of Green's question has a negative answer. Vanishing of
$\mu$ does not mean $P=1$: in finite $D_4$, with $x$ the product of the
three leaves, $b=xcx$ has maximum common commuting descents and even gap
$4$, so $\mu(x,b)=0$, but $P_{x,b}=1+2q$
\cite[Corollary~3 (arXiv v1, \S4)]{Mongelli}.

\subsection{Terminals in affine \texorpdfstring{$E_8$}{E8} beyond the family}
\label{sec:affineextra}
The reflection
\[
 w=r_\beta=875634523412301282341230128567453,\qquad
 \beta=(1,2,3,3,2,2,1,1,2),
\]
of length $33$ is a full-support terminal with $L(w)=R(w)=\{3,5,7,8\}$,
which is not a maximum independent set, so Proposition~\ref{prop:parity}
does not apply. It has exactly five eligible fully commutative lower
endpoints, all with $\deg P_{x,w}=13$:
\begin{center}\small
\begin{tabular}{@{}lrrrrr@{}}
\toprule
$x$ & $8753$ & $87530$ & $87531$ & $875301$ & $875310$\\
gap $\ell(w)-\ell(x)$ & $29$ & $28$ & $28$ & $27$ & $27$\\
$\mu(x,w)$ & $0$ & $0$ & $0$ & $1$ & $1$\\
\bottomrule
\end{tabular}
\end{center}
The right extensions $ws_0$ and $ws_1$ of length $34$ are terminal
involutions with $L=R=\{0,3,5,7,8\}$, respectively $\{1,3,5,7,8\}$, each a
maximum independent set, each with a unique eligible fully commutative
lower endpoint ($87530$, respectively $87531$, gap $29$, $\deg P=14$) and
$\mu=1$, so these two pairs satisfy every hypothesis of
Proposition~\ref{prop:parity} except parity. All values found lie in
$\{0,1\}$. The polynomials and the sizes of the lower ideals and
quotients are recorded in the certificate
(Appendix~\ref{app:certificates}); the computation runs the
Kazhdan--Lusztig recursion on the parabolic quotient $W/W_{R(w)}$, which
is legitimate because $P_{x,w}=P_{xs,w}$ for $s\in R(w)$
\cite[\S2.3]{KL}.

\subsection{Conditional interval transfer}
Combinatorial invariance, conjectured independently by Lusztig and Dyer
\cite{Dyer}, asserts that $P_{x,w}$ depends only on the poset $[x,w]$; it
is known for lower intervals \cite{Brenti04,BCM,Delanoy}, for the
coefficient of $q$ \cite{Patimo,BGL} and for short intervals
\cite{Incitti,BGL,EMS} (see also \cite{BBDVW}), and a complete proof for
arbitrary Coxeter systems is claimed in the unrefereed machine-generated
OpenAI preprint \cite[Theorem~1.1]{CI}. Assuming it, the bound of Corollary~\ref{cor:finiteADE} and the
vanishing for the families of Section~\ref{sec:families} hold for every
Bruhat interval in any Coxeter group that is isomorphic as a poset to one of
the intervals concerned; the formal conditional consequences for Gern's
intervals $[x_n,w_n]$ in $D_n$ are Theorem~2.1 and Corollary~2.2 of the
archived note \texttt{kl-transfer} (Appendix~\ref{app:certificates}). No new
interval isomorphisms are asserted.

\subsection{Open problems}
(1) Is $\mu(x,w)\in\{0,1\}$ for every fully commutative $x$ and every
$w\in W(\widetilde E_8)$? All values computed so far, including those of
Section~\ref{sec:affineextra}, lie in $\{0,1\}$.
(2) Classify the terminal elements of $\widetilde E_8$, and of the
indefinite groups $E_n$, up to translation.
(3) Determine $\mu$ for the translated length-$33$ family
$r_{\beta+k\delta'}$, $k\geq1$.
(4) Settle Green's question for $F_4$, $H_3$ and $H_4$.
(5) Give an effective description of all noncommuting terminals of $E_n$
for $n\geq9$; Theorem~\ref{thm:uniform} shows that no finite list exists.

\section*{Use of AI tools}
This work was carried out with substantial assistance from large language
model systems. OpenAI's GPT-6.1 Sol and GPT-6 Astra, run as teams of
agents, performed the literature searches, proposed and drafted the
constructions and proofs, wrote the enumeration and verification code,
produced internal review notes, and drafted the text. Anthropic's Claude
Fable 5.1 then carried out an independent adversarial review: it
re-derived every lemma against the primary sources, recomputed every
computational claim with separately written code, checked every citation,
and found the strengthening of Theorem~\ref{thm:uniform}, the data in
Section~\ref{sec:complements}, and the attribution corrections incorporated
here; the Go port of the supplement was written by Claude Sonnet 5.5 and
verified by Claude Fable 5.1 against the earlier programs' outputs. The
author directed the work, checked the proofs, computations and
references, and takes full responsibility for the content. The complete
workflow, including the prompts used, the internal review notes (which are
machine-generated and are not peer review), and the review code, is
recorded in the repository cited in Appendix~\ref{app:certificates}.

\appendix
\section{Computational certificates}\label{app:certificates}
\begin{table}[t]
\centering\small
\begin{tabular}{@{}>{\raggedright\arraybackslash}p{0.17\textwidth}>{\raggedright\arraybackslash}p{0.80\textwidth}@{}}
\toprule
Result & Programs; certificates\\
\midrule
$E_6$, $E_7$ & \texttt{enumerate-bad}, \texttt{matrix-search}, \texttt{verify-exceptional};
 \texttt{e6\_bad.json}, \texttt{e7\_bad.json}, \texttt{e6-independent-certificate.json}, \texttt{e7-independent-certificate.json}\\
$E_8$ & \texttt{terminals-flat}, \texttt{terminals-recursive} (both on \texttt{internal/parabolic}),
 \texttt{verify-outputs}, \texttt{e8-d7}, \texttt{verify-e8};
 \texttt{e8-recursive.json}$^\dagger$, \texttt{e8-d7-terminals.json}$^\dagger$\\
$E_8$, four chains & \texttt{terminals-recursive -chains-certificate} ($E_7$, $D_7$, $A_7$ and a second $E_7$ chain); \texttt{e8-recursive-chains.json}\\
Table~\ref{tab:terminals}, Remark~\ref{rem:gern} & \texttt{finite-descents}, \texttt{terminal-structure}$^*$;
 \texttt{finite-descents.json}, \texttt{terminal-structure-certificate.json}\\
Table~\ref{tab:terminals}, third check & \texttt{terminal-data-check}; \texttt{checks.json}\\
$D_6$, $D_8$ & \texttt{d6-certificate} (on \texttt{internal/d6}),
 \texttt{d8-gern-kl}$^*$; \texttt{d6-recurrence-certificate.json}, \texttt{d8-gern-certificate.json}\\
$D_6$ inside $E_7$, $E_8$ & \texttt{d6-ambient}$^*$; \texttt{d6-ambient-certificate.json}\\
Uniform family & \texttt{uniform-verify}, \texttt{uniform-covers}$^*$,
 \texttt{fc-maxima}$^*$; \texttt{referee-certificate.json}$^\dagger$, \texttt{uniform-family-certificate.json}, \texttt{fc-maxima-certificate.json}\\
Affine $\widetilde E_8$ & \texttt{affine-proof}, \texttt{affine-fc-covers} (both on \texttt{internal/affine}),
 \texttt{fc-catalogue};
 \texttt{proved-affine-certificate.json}, \texttt{fc-cover-certificate.json}$^\dagger$\\
$E_{10}$ & \texttt{e10-all-k} (on \texttt{internal/indefinite}); \texttt{e10-certificate.json}$^\dagger$\\
Section~\ref{sec:complements} & \texttt{e6-mu-table}$^*$, \texttt{e9-quotient-kl}$^*$,
 \texttt{affine-d4};
 \texttt{e6-mu-table.json}, \texttt{e9-odd-gap-certificate.json}, \texttt{affine-d4-independent.json}\\
Section~\ref{sec:affineextra} & \texttt{affine-reflection-family}; \texttt{affine\_reflection\_family.json}\\
\bottomrule
\end{tabular}
\caption{Programs (the commands \texttt{cmd/\textit{name}} of the Go module)
and certificates of the supplement (the zip of release v0.4.0), by result.
Starred programs were written for the review and share no code with the
others; the \texttt{internal} libraries in parentheses are the only shared
code; daggered certificates are written by the runner, not shipped as
inputs.}
\label{tab:programs}
\end{table}
The following computations are parts of proofs: the terminal
classification of Proposition~\ref{prop:classification}; the reducedness,
involution and descent data of the words in Table~\ref{tab:terminals} and
of $c_0$, whose matrix (hence the columns $c_0(\alpha_j)$) is recorded; the
$D_6$ polynomial
$P_{x_6,w_6}=1+6q+11q^2+6q^3+q^4+q^5$; the complete fully commutative
catalogue of $\widetilde E_8$ ($44{,}199$ elements, maximum length $44$)
and the cover check for $c_0$; and the lowering words certifying the real
roots $\gamma$, $\gamma+\delta$, $\hat\beta_0$, $\hat\gamma$ and
$\hat\gamma+\hat\delta$. Everything else is a cross-check or a complement:
the fully commutative catalogues of $E_6$, $E_7$, $E_8$ (agreeing with
\cite{Stembridge98} and \cite[\S5.1]{BJN}), Poincar\'e-polynomial checks,
finite-rank checks of the uniform family, the structure data of
Remark~\ref{rem:gern}, the fully commutative catalogues and maxima of
$E_{10}$--$E_{13}$, the data of Section~\ref{sec:complements}, and the
$D_8$ polynomial
$P_{x_8,w_8}=P_{e,w_8}=1+12q+59q^2+154q^3+233q^4+221q^5+147q^6+70q^7
+20q^8+2q^9$ of the archived note \texttt{kl-transfer}. All arithmetic is
exact. For $E_6$ and $E_7$ two programs sharing no code (root indices
versus integer columns) give full enumerations agreeing with the
Poincar\'e products; for $E_8$ the recursive pruning engine and a separate
enumeration along a $D_7$ parabolic reproduce the same list, sharing only
their packed root representation, an integer-matrix verifier rechecks
every terminal word, length, terminal test and eligible lower endpoint,
and \texttt{terminals-recursive} gives the same list along four parabolic
chains ($E_7$, $D_7$, $A_7$ and a second $E_7$ chain);
\texttt{terminal-data-check} rechecks the data of Table~\ref{tab:terminals}
from the certificates with separate code. The $D_6$
polynomial comes from a sparse Kazhdan--Lusztig recursion checked through
all $24{,}245$ dependency records of the saved certificate (a consistency
check, the checker sharing group arithmetic with the generator);
\texttt{d8-gern-kl} recomputes it with independent code in the
signed-permutation model and, in the extended run, in the geometric
representation as well, and \texttt{d6-ambient} recomputes it directly
inside $E_7$, $E_8$ and the $D_6$ parabolic ($3{,}184$-element lower
ideal, geometric representation). The uniform root proof is
symbolic; a verifier checks
Lemmas~\ref{lem:uniformseed}--\ref{lem:uniformdescents} symbolically and
the finite-rank matrices for $3\leq r\leq40$. The starred programs of
Table~\ref{tab:programs}, written for the review, share no code with the
others.

Programs, inputs and output certificates (Table~\ref{tab:programs}) are
archived with SHA-256 checksums in release v0.4.0 of the repository
\texttt{github.com/tygern/kl} (zip \texttt{exceptional-leading-proof.zip});
a single command rebuilds every classification and verifies every
certificate with Go 1.22 or later and no other dependency. The earlier
Python and C++ programs remain in release v0.2.0, and the Go port was
verified against their outputs certificate for certificate. Copies of the daggered
certificates of Table~\ref{tab:programs}, which the runner writes rather
than reads, are kept in the repository. The archived research notes of
Section~\ref{sec:complements} are in \texttt{research/archive/notes/}.

\begin{thebibliography}{99}\footnotesize\raggedright
\setlength{\itemsep}{0pt}\setlength{\parsep}{0pt}\setlength{\parskip}{0pt}
\bibitem{BGL} G. T. Barkley, C. Gaetz, and T. Lam, Combinatorial
invariance for the coefficient of $q$ in Kazhdan--Lusztig polynomials,
preprint (2026), arXiv:2601.07793.
\bibitem{BJN} R. Biagioli, F. Jouhet, and P. Nadeau, Fully commutative
elements in finite and affine Coxeter groups, Monatsh. Math. 178 (2015),
1--37, doi:10.1007/s00605-014-0674-7.
\bibitem{BB} A. Bj\"orner and F. Brenti, Combinatorics of Coxeter groups,
Grad. Texts in Math. 231, Springer, 2005, doi:10.1007/3-540-27596-7.
\bibitem{BBDVW} C. Blundell, L. Buesing, A. Davies, P. Veli\v{c}kovi\'c,
and G. Williamson, Towards combinatorial invariance for Kazhdan--Lusztig
polynomials, Represent. Theory 26 (2022), 1145--1191, doi:10.1090/ert/624.
\bibitem{Brenti04} F. Brenti, The intersection cohomology of Schubert
varieties is a combinatorial invariant, European J. Combin. 25 (2004),
1151--1167, doi:10.1016/j.ejc.2003.10.011.
\bibitem{BCM} F. Brenti, F. Caselli, and M. Marietti, Special matchings
and Kazhdan--Lusztig polynomials, Adv. Math. 202 (2006), 555--601,
doi:10.1016/j.aim.2005.01.011.
\bibitem{Chmutov} M. S. Chmutov, The structure of $W$-graphs arising in
Kazhdan--Lusztig theory, Ph.D. thesis, Univ. of Michigan, 2014,
hdl:2027.42/108802.
\bibitem{Delanoy} E. Delanoy, Combinatorial invariance of Kazhdan--Lusztig
polynomials on intervals starting from the identity, J. Algebraic Combin.
24 (2006), 437--463, doi:10.1007/s10801-006-0014-7.
\bibitem{Dyer} M. Dyer, On the ``Bruhat graph'' of a Coxeter system,
Compositio Math. 78 (1991), 185--191,
\url{https://www.numdam.org/item/CM_1991__78_2_185_0/}.
\bibitem{Ernst} D. C. Ernst, Non-cancellable elements in type affine $C$
Coxeter groups, Int. Electron. J. Algebra 8 (2010), 191--218,
arXiv:0910.0923.
\bibitem{EMS} F. Esposito, M. Marietti, and S. Stella, Flip combinatorial
invariance and Weyl groups, preprint (2025; v2 2026), arXiv:2509.16433.
\bibitem{Fan} C. K. Fan, Structure of a Hecke algebra quotient,
J. Amer. Math. Soc. 10 (1997), 139--167, doi:10.1090/S0894-0347-97-00222-1.
\bibitem{FanGreen} C. K. Fan and R. M. Green, On the affine
Temperley--Lieb algebras, J. London Math. Soc. (2) 60 (1999), 366--380,
doi:10.1112/S0024610799007796.
\bibitem{GeckPyCox} M. Geck, PyCox: Computing with (finite) Coxeter groups
and Iwahori--Hecke algebras, LMS J. Comput. Math. 15 (2012), 231--256,
doi:10.1112/S1461157012001064.
\bibitem{GeckHalls} M. Geck and A. Halls, On the Kazhdan--Lusztig cells in
type $E_8$, Math. Comp. 84 (2015), 3029--3049, doi:10.1090/mcom/2963.
\bibitem{Gern} T. C. Gern, Leading coefficients of Kazhdan--Lusztig
polynomials in type $D$, Ph.D. thesis, Univ. of Colorado Boulder, 2013,
arXiv:1304.6074.
\bibitem{Graham} J. J. Graham, Modular representations of Hecke algebras
and related algebras, Ph.D. thesis, Univ. of Sydney, 1995.
\bibitem{GreenStar} R. M. Green, Star reducible Coxeter groups, Glasgow
Math. J. 48 (2006), 583--609, doi:10.1017/S0017089506003211.
\bibitem{GreenJones} R. M. Green, Generalized Jones traces and
Kazhdan--Lusztig bases, J. Pure Appl. Algebra 211 (2007), 744--772,
doi:10.1016/j.jpaa.2007.04.004.
\bibitem{GreenEn} R. M. Green, On the Markov trace for Temperley--Lieb
algebras of type $E_n$, J. Knot Theory Ramifications 18 (2009), 237--264,
doi:10.1142/S0218216509006872, arXiv:0704.0283 (statement numbers cited
here follow v1).
\bibitem{GreenLeading} R. M. Green, Leading coefficients of
Kazhdan--Lusztig polynomials and fully commutative elements, J. Algebraic
Combin. 30 (2009), 165--171, doi:10.1007/s10801-008-0156-x,
arXiv:0801.1650.
\bibitem{GreenLosonczy} R. M. Green and J. Losonczy, Fully commutative
Kazhdan--Lusztig cells, Ann. Inst. Fourier 51 (2001), 1025--1045,
doi:10.5802/aif.1843.
\bibitem{Incitti} F. Incitti, More on the combinatorial invariance of
Kazhdan--Lusztig polynomials, J. Combin. Theory Ser. A 114 (2007),
461--482, doi:10.1016/j.jcta.2006.06.009.
\bibitem{Jones} B. C. Jones, Leading coefficients of Kazhdan--Lusztig
polynomials for Deodhar elements, J. Algebraic Combin. 29 (2009), 229--260,
doi:10.1007/s10801-008-0131-6.
\bibitem{KL} D. Kazhdan and G. Lusztig, Representations of Coxeter groups
and Hecke algebras, Invent. Math. 53 (1979), 165--184,
doi:10.1007/BF01390031.
\bibitem{Laird} T. Laird, A study of $T$-avoiding elements of Coxeter
groups, M.S. thesis, Northern Arizona Univ., 2016,
\url{https://in.nau.edu/department-mathematics-statistics/wp-content/uploads/sites/377/2018/08/LairdThesis-ek.pdf}.
\bibitem{LusztigII} G. Lusztig, Cells in affine Weyl groups, II,
J. Algebra 109 (1987), 536--548, doi:10.1016/0021-8693(87)90154-2.
\bibitem{LusztigUP} G. Lusztig, Hecke algebras with unequal parameters,
CRM Monogr. Ser. 18, Amer. Math. Soc., 2003, doi:10.1090/crmm/018.
\bibitem{Mongelli} P. Mongelli, Kazhdan--Lusztig polynomials of boolean
elements, J. Algebraic Combin. 39 (2014), 497--525,
doi:10.1007/s10801-013-0455-8, arXiv:1111.2945.
\bibitem{CI} OpenAI, Combinatorial invariance of Kazhdan--Lusztig
polynomials, OpenAI Math Release preprint, 24 September 2026 (not peer
reviewed), commit adc7f124, SHA-256
7e1be2107a85e28d9ec339b4ee62d4efeb909968df6768368e2779710624e331,
\url{https://github.com/openai/math/blob/main/preprints/Combinatorial-Invariance-of-Kazhdan-Lusztig-Polynomials-September-24-2026/paper.pdf}.
\bibitem{Patimo} L. Patimo, A combinatorial formula for the coefficient of
$q$ in Kazhdan--Lusztig polynomials, Int. Math. Res. Not. IMRN 2021,
no.~5, 3203--3223, doi:10.1093/imrn/rnz255.
\bibitem{Shi} J.-Y. Shi, Fully commutative elements in the Weyl and affine
Weyl groups, J. Algebra 284 (2005), 13--36,
doi:10.1016/j.jalgebra.2004.08.036.
\bibitem{Stembridge} J. R. Stembridge, On the fully commutative elements
of Coxeter groups, J. Algebraic Combin. 5 (1996), 353--385,
doi:10.1007/BF00193185; author preprint
\url{https://websites.umich.edu/~jrs/papers/FC.pdf} (its Proposition~1.1 and
Theorem~4.1 are Proposition~2.1 and Theorem~5.1 of the published version).
\bibitem{Stembridge98} J. R. Stembridge, The enumeration of fully
commutative elements of Coxeter groups, J. Algebraic Combin. 7 (1998),
291--320, doi:10.1023/A:1008623323374.
\end{thebibliography}
\end{document}
